\section{Introduction}
\label{sec:intro}

Large multimodal models inherit their safety alignment from a language-model backbone whose RLHF training was performed entirely on text. Every additional input or output modality grafted onto that backbone---images, audio, video, or, in our case, 7-DoF robot action tokens---creates a pathway that the refusal circuits have never been shown. When the alignment transfers, we get safer multimodal assistants; when it does not, at best we get multimodal hallucinations, and at worst, physical harm.

Vision-Language-Action (VLA) models sit at the most consequential end of this spectrum. An unsafe action is not an unsafe \emph{sentence}: it is a gripper closing on a child's hand, a liquid poured into an open flame, a weapon aimed at a person. Recent adversarial work shows that every open VLA studied so far can be induced to emit such actions through plain-language instructions or small image perturbations \cite{wu2024exploringadversarial,badrobot2025,shawshank2025}, and existing VLA defenses either retrain the entire policy end-to-end in a safe-RL loop \cite{safevla2025} or rely on geometric collision avoidance that carries no semantic notion of harm \cite{vlsa2025}.

\paragraph{A structural safety gap in action-token VLAs.}
We observe that \textsc{OpenVLA} \cite{openvla2024}, the de-facto open VLA, is built on Llama-2-7b-\emph{base}---a model that was \emph{never RLHF-aligned}---and encodes its action vocabulary by overwriting the last 256 tokens of the Llama vocabulary with learned discrete bins. These two choices are standard practice; we show that, in combination, they create a geometric gap that existing multimodal safety defenses cannot cross. First, the backbone has no refusal direction of its own: it was never taught to refuse. Second, the action tokens, after extensive behavior-cloning fine-tuning, live in a residual-stream subspace that is nearly orthogonal to where the natural-language ``harmful/harmless'' axis would have been had the backbone been chat-aligned. We quantify this gap empirically in Sec.~\ref{sec:method:diagnostic}: a safety-layer probe that identifies a clear harmful/benign separation in aligned LVLMs \cite{securitytensors2025} produces a markedly weaker signal in \textsc{OpenVLA}'s action-decoding path.

\paragraph{Cross-model refusal direction transplant.}
We propose a fix that does not require retraining the VLA. \methodfull{} (\method{}) extracts a refusal direction from Llama-2-7b-\emph{chat}---a sibling of \textsc{OpenVLA}'s backbone that shares its pretraining initialization but did receive RLHF alignment---and re-injects it at inference time only at action-token positions of \textsc{OpenVLA}. The method is a single forward hook and composes with existing prompt-level defenses.

Two prior works are adjacent but insufficient for our setting. \textsc{VLM-Guard} \cite{vlmguard2025} extracts a similar steering direction but applies it to every token, over-perturbing benign vision-language processing. \textsc{SARSteer} \cite{sarsteer2025} steers text-derived refusal directions in audio LMs, but crucially extracts the direction from the target model itself, which is zero-shot inapplicable when the target backbone (Llama-2-base) has no refusal responses to extract.

\paragraph{Contributions.}
\begin{itemize}
    \item We identify a structural safety gap specific to action-token VLAs, arising from the combination of an \emph{unaligned} backbone and a learned discrete action vocabulary (Sec.~\ref{sec:method:diagnostic}).
    \item We introduce \method{} and \method{}+, training-free inference-time methods that transplant refusal directions from an externally-aligned LLM into a VLA. Base \method{} targets action-token positions; \method{}+ adds a generation-phase-aware dual injection covering all 7 action tokens including $a_1$ (Sec.~\ref{sec:method:rdt}--\ref{sec:method:dualphase}), together with a rank-$k$ subspace variant (Sec.~\ref{sec:method:rankk}) and a content-adaptive scale head (Sec.~\ref{sec:method:adaptive}).
    \item We design an evaluation across three suites (\textsc{LIBERO-Long}, \textsc{SafeAgentBench}, an instruction-conditioned \textsc{AdvBench} split) and four attack families (textual jailbreak, \textsc{UADA}, \textsc{UPA}, \textsc{TMA} patches), comparing \method{}+ against prompt-level, hidden-state-\allowbreak projection, and LoRA fine-tuning baselines (Sec.~\ref{sec:exp}).
    \item We include a direction-source control experiment---refusal direction vs.\ random unit vector of the same norm---to confirm that the cross-model geometry of the refusal direction, not merely hidden-state perturbation, is the active ingredient (Sec.~\ref{sec:exp:ablation}).
    \item We stress-test \method{}+ against a white-box adaptive attacker and analyze what ``refusal'' means in action space, including a safety-layer probe recovery analysis (Sec.~\ref{sec:exp:adaptive}, Sec.~\ref{sec:analysis:action-space}, Sec.~\ref{sec:analysis:probe-recovery}).
\end{itemize}
